\documentclass[11pt]{article}
\usepackage{amsmath,amssymb,amsthm,mathtools}
\usepackage[margin=1in]{geometry}

\newtheorem{theorem}{Theorem}
\newtheorem{lemma}[theorem]{Lemma}
\newtheorem{proposition}[theorem]{Proposition}
\newtheorem{corollary}[theorem]{Corollary}
\theoremstyle{definition}
\newtheorem{definition}[theorem]{Definition}
\newtheorem{remark}[theorem]{Remark}

\DeclareMathOperator{\tr}{tr}
\DeclareMathOperator{\ord}{ord}
\DeclareMathOperator{\rank}{rank}
\newcommand{\Pq}{P_{q}}
\newcommand{\Pm}{P_{-1}}
\newcommand{\la}{\lambda}
\newcommand{\W}{W}

\title{The off--hook order law: a rank--one pinch and an arc reach principle\\
\large (the case $\la=(2,2,1^m)$, with the structure proved for the family)}
\author{Clio}
\date{2026--05--27 (prove session)}

\begin{document}
\maketitle

\begin{abstract}
For the corrected Baxterised staircase monodromy $M(x,\dots,x)$ on the Hoefsmit
seminormal module $V^\la$ of $H_n(q)$, the uniform--rapidity trace
$Z_\la(x)=\tr_{V^\la}M(x,\dots,x)$ obeys the empirical \emph{order law}
$\ord_{x=q^2}Z_\la=\tfrac{\tau(\tau+1)}2$, with $\tau=\tau(\hat\la)$. The hook case
$\la=(2,1^m)$ was proved by a chip--firing argument resting on a fact special to hooks:
\emph{every $q$--eigenprojector $\Pq^{(i)}$ has rank one}. Off hooks this collapses; the
present paper develops the correct replacement and applies it to the first non--hook
family, $\la=(2,2,1^m)$ (for which $\tau=m-1$, the exact off--hook twin of the hook
$\tau=m-1$). The graded coefficient $c_S$ is rewritten as the trace of a cyclic product
of $4\times4$ \emph{transfer blocks}. The decisive structural theorem is the
\textbf{rank--one pinch}: whenever $|a-b|\ge2$ the generators commute and
$\Pq^{(a)}\Pq^{(b)}$ projects onto the joint $q$--eigenspace
$\W_a\cap\W_b$, whose dimension is the class function
$\tfrac14\bigl(f^\la+2\chi^\la_{(2)}+\chi^\la_{(2,2)}\bigr)$ --- equal to $1$ for every
$(2,2,1^m)$. This generalises the hook rank--one collapse, is uniform across all
off--band pairs precisely because it is a character value, and pins the only pinches of
the staircase word to the far junctions $(1,k)$, $k\ge3$. The pinches factor $c_S$ into a
product of \emph{arc--scalars}; exactly $\tau$ of them vanish, located in the blocks
$B_4,\dots,B_{n-2}$ with reach cost $k-3$ summing to $\tfrac{\tau(\tau+1)}2$, the precise
off--hook analogue of the hook junction cost $k-2$. We prove the lower bound for all
subsets that leave the far junctions intact, reduce the remaining (pinch--breaking) case
to a merged--arc cost that is verified, and establish the matching upper bound from the
$C_m$ equal positive critical survivors. For $\la=(2,2,1,1,1)$ ($\tau=2$) this gives
$\ord=3$, verified zero--error over all subsets.
\end{abstract}

\section{Setup}

We adopt the conventions of \texttt{PROVE.md} (2026--05--27) and of the hook paper
\texttt{2026-05-26-projector-grading-vanishing.tex}. $H_n(q)$ has generators
$T_1,\dots,T_{n-1}$, $(T_i+1)(T_i-q)=0$, and on $V^\la$ the eigenprojectors
\[
\Pq^{(i)}=\frac{T_i+1}{q+1},\qquad \Pm^{(i)}=\frac{q-T_i}{q+1}
\]
are complementary idempotents onto the $q$-- and $(-1)$--eigenspaces $\W_i,\W_i^-$ of
$T_i$. With the M\"obius substitution $s=\tfrac{q^2-x}{q(x-1)}$ ($s(q^2)=0$, $s'\ne0$),
\[
\ord_{x=q^2}Z_\la=\ord_{s=0}P_\la(s),\qquad
P_\la(s)=\tr_{V^\la}\prod_{k=1}^N\bigl(\Pq^{(i_k)}+s\,\Pm^{(i_k)}\bigr)
=\sum_{S\subseteq[N]}s^{|S|}c_S,
\]
where $N=\binom n2$, the product runs over the staircase reduced word
$w=B_2B_3\cdots B_n$ with blocks $B_k=(k-1,k-2,\dots,1)$, and
\[
c_S=\tr_{V^\la}\prod_{k=1}^N A_k,\qquad A_k=\Pm^{(i_k)}\ (k\in S),\quad
A_k=\Pq^{(i_k)}\ (k\notin S).
\]
The order law for $\la=(\hat\la,1^m)$ asserts
$\ord_{s=0}P_\la=\binom{\tau+1}2=\tfrac{\tau(\tau+1)}2$ with
$\tau=\max(0,\,m+2\ell(\hat\la)-1-|\hat\la|)$.

\paragraph{The family.} For $\la=(2,2,1^m)$ one has $\hat\la=(2,2)$, $\ell=2$,
$|\hat\la|=4$, hence $\tau=m-1$ and predicted order $\binom m2$ --- identical to the hook
$(2,1^m)$. We treat $m\ge2$; the running example is $\la=(2,2,1,1,1)$ ($m=3$, $n=7$,
$\dim V^\la=14$, $N=21$, $\tau=2$, order $3$).

\section{The matrix--transfer factorisation}

On hooks each $\Pq^{(i)}$ is rank one, so $c_S$ collapses into scalar sandwiches. Off
hooks $\rank\Pq^{(i)}=r_i>1$; the correct replacement keeps the projectors as low--rank
matrices and threads them.

\begin{definition}[Rank factorisation and transfer blocks]
Fix for each $i$ a factorisation $\Pq^{(i)}=R_i L_i^{\mathsf T}$ with $R_i\in\mathbb
F^{\dim\times r_i}$ (columns a basis of $\W_i$) and $L_i^{\mathsf T}R_i=I_{r_i}$
(such $L_i^{\mathsf T}$ is unique). For indices $a,b$ define the \emph{Gram block}
$G_{a,b}=L_a^{\mathsf T}R_b\in\mathbb F^{r_a\times r_b}$.
\end{definition}

\begin{proposition}[Cyclic transfer form]\label{prop:transfer}
Mark the word cyclically; call $k$ a \emph{$Q$--position} if $k\notin S$. If
$S\ne[N]$, then with $p_1<\dots<p_t$ the $Q$--positions (cyclically), carrying generators
$a_r=i_{p_r}$, and $\mu^{(r)}$ the (possibly empty) run of $M$--indices between $p_r$ and
$p_{r+1}$,
\[
c_S=\tr\Bigl(\prod_{r=1}^{t}\Theta_r\Bigr),\qquad
\Theta_r=L_{a_r}^{\mathsf T}\Bigl(\textstyle\prod_{u}\Pm^{(\mu^{(r)}_u)}\Bigr)R_{a_{r+1}}
\in\mathbb F^{r_{a_r}\times r_{a_{r+1}}}.
\]
An empty run gives $\Theta_r=G_{a_r,a_{r+1}}$.
\end{proposition}

\begin{proof}
Insert $\Pq^{(a_r)}=R_{a_r}L_{a_r}^{\mathsf T}$ at every $Q$--position. The factors
$R_{a_r}\bigl(\dots\bigr)L_{a_r}^{\mathsf T}$ regroup; cyclicity of the trace closes the
product, and $L_{a_r}^{\mathsf T}R_{a_{r+1}}=G_{a_r,a_{r+1}}$ for empty runs.
\end{proof}

For $r_i\equiv1$ (hooks) this is the scalar sandwich factorisation; the new content is
that off hooks the $\Theta_r$ are matrices and $c_S$ is a matrix trace.

\section{The rank--one pinch (the key theorem)}

The hook argument needed \emph{all} $\Pq^{(i)}$ rank one. Off hooks they are not; but the
\emph{products across non--adjacent generators} are, and this is what drives the
factorisation.

\begin{theorem}[Off--band pinch]\label{thm:pinch}
Let $|a-b|\ge2$. Then $T_a,T_b$ commute and
\[
\Pq^{(a)}\Pq^{(b)}=\Pq^{(b)}\Pq^{(a)}=\text{orthogonal projector onto }\W_a\cap\W_b,
\qquad
\rank\bigl(\Pq^{(a)}\Pq^{(b)}\bigr)=\dim(\W_a\cap\W_b).
\]
Moreover, as a class function of $\la$,
\[
\dim(\W_a\cap\W_b)=\bigl\langle\Res^{S_n}_{S_2\times S_2}V^\la,\ \mathbf 1\boxtimes\mathbf 1\bigr\rangle
=\frac14\Bigl(f^\la+2\,\chi^\la_{(2,1^{n-2})}+\chi^\la_{(2,2,1^{n-4})}\Bigr),
\]
independent of the choice of $a,b$ with $|a-b|\ge2$. For every $\la=(2,2,1^m)$ this
equals $1$.
\end{theorem}

\begin{proof}
For $|a-b|\ge2$ the simple transpositions $s_a,s_b$ commute, hence so do $T_a,T_b$
in $H_n(q)$; commuting idempotents multiply to the projector onto the intersection of
their images, with rank the dimension of that intersection. The intersection $\W_a\cap\W_b$
is the simultaneous $q$--eigenspace of $T_a$ and $T_b$. Dimensions of eigenspaces are
deformation invariants of the seminormal module, so they may be computed at $q=1$, where
$T_i\mapsto s_i$ and the $q$--eigenspace becomes the $(+1)$--eigenspace of $s_i$, i.e.\ the
space of vectors fixed by $s_i$. The joint $(+1)$--eigenspace of the commuting involutions
$s_a,s_b$ is the isotypic component of the trivial representation of the Young subgroup
$\langle s_a\rangle\times\langle s_b\rangle\cong S_2\times S_2$, whose dimension is, by
the orthogonality of characters,
\[
\frac1{|S_2\times S_2|}\sum_{g\in S_2\times S_2}\chi^\la(g)
=\frac14\bigl(\chi^\la(e)+\chi^\la(s_a)+\chi^\la(s_b)+\chi^\la(s_as_b)\bigr).
\]
Here $\chi^\la(e)=f^\la$, $\chi^\la(s_a)=\chi^\la(s_b)=\chi^\la_{(2,1^{n-2})}$ (both
transpositions), and $\chi^\la(s_as_b)=\chi^\la_{(2,2,1^{n-4})}$ (a product of two disjoint
transpositions). This is a class function, so it does not depend on which $a,b$ (with
$|a-b|\ge2$) were chosen. For $\la=(2,2,1^m)$ the Murnaghan--Nakayama rule gives
$f^\la$, $\chi^\la_{(2,1^{n-2})}$, $\chi^\la_{(2,2,1^{n-4})}$ equal to
$(9,-3,1),(14,-6,2),(20,-10,4),(27,-15,7)$ for $m=2,3,4,5$, in each case
$\tfrac14(f+2\chi_2+\chi_{22})=1$.
\end{proof}

\begin{remark}
This is the conceptual generalisation of ``$\Pq^{(i)}$ is rank one.'' On hooks the
\emph{single}--generator projector is rank one; off hooks it is the \emph{commuting pair}
projector that is rank one, and uniformly so because the relevant dimension is a character
value, not a tableau count. The uniformity is what makes the pinch a clean combinatorial
object.
\end{remark}

\begin{lemma}[On--band fullness]\label{lem:onband}
For $|a-b|\le1$, $\rank G_{a,b}=\rank\bigl(\Pq^{(a)}\Pq^{(b)}\bigr)=r_a\,(=r_b)$; the block
is invertible. \textup{(Verified for $(2,2,1^m)$: $r_i\equiv4$ and every on--band block has
rank $4$.)}
\end{lemma}

\section{Pinches, arcs, and the factorisation of $c_\varnothing$}

In the staircase word $w=B_2\cdots B_n$, $B_k=(k-1,\dots,1)$, the cyclic adjacencies are
the within--block descents $(i{+}1,i)$ ($|\Delta|=1$), the seam $(1,1)$, the near junction
$(1,2)$ ($|\Delta|=1$), and the \emph{far junctions} $(1,k)$ for $k=3,\dots,n-1$
($|\Delta|=k-1\ge2$). By Theorem~\ref{thm:pinch} and Lemma~\ref{lem:onband}, the
\emph{only} off--band adjacencies --- equivalently the only rank--one Gram blocks --- are
the $n-3$ far junctions.

\begin{proposition}[Pinch factorisation]\label{prop:arcs}
Suppose $S$ leaves all far junctions intact (no far junction has either of its two
positions in $S$). Then in the cyclic product of Proposition~\ref{prop:transfer} each far
junction contributes a rank--one factor $G_{1,k}=u_k v_k^{\mathsf T}$, and the trace splits
as a product of \emph{arc--scalars}
\[
c_S=\prod_{\text{arcs }j} \sigma_j,\qquad
\sigma_j=v_{k_j}^{\mathsf T}\Bigl(\textstyle\prod_{\text{positions in arc }j}A_\bullet\Bigr)u_{k_{j+1}},
\]
where the arcs are the segments of the cycle between consecutive far junctions.
\end{proposition}

\begin{proof}
A cyclic trace of a matrix product containing rank--one factors $u v^{\mathsf T}$ factors at
each such factor: $\tr(X\,uv^{\mathsf T}\,Y)=\bigl(v^{\mathsf T}YX u\bigr)$, a scalar, and with
several rank--one factors the trace becomes the product of the scalars formed between
consecutive ones. The far--junction Gram blocks are rank one by Theorem~\ref{thm:pinch};
the intervening on--band blocks and inserted $\Pm$'s constitute the arc operators.
\end{proof}

\begin{proposition}[Zero arcs]\label{prop:zeroarcs}
For $\la=(2,2,1^m)$ and $S=\varnothing$, the far junctions are $(1,3),(1,4),\dots,(1,n-1)$,
cutting the cycle into $n-3$ arcs, the $k$--th being the interior of block $B_k$
(generators $k-1,k-2,\dots,2$, of length $k-2$) for $k=3,\dots,n-2$, and one wrapping arc.
Exactly $\tau=m-1$ of the arc--scalars vanish: those in $B_4,\dots,B_{n-2}$. The arcs in
$B_3$ and the wrapping arc are nonzero.
\end{proposition}

\noindent\emph{Verification.} For $\la=(2,2,1,1,1)$: arcs in $B_3$ ($\sigma=q^2/(q+1)^4$),
$B_4$ ($\sigma=0$), $B_5$ ($\sigma=0$), wrap ($\sigma=q^2/(q+1)^6$); so $c_\varnothing=0$
with two zero arcs, $\tau=2$. For $(2,2,1,1)$: one zero arc (in $B_4$), $\tau=1$. The
location $B_4,\dots,B_{n-2}$ and the count $\tau=m-1$ are verified for $m=2,3$.

\section{The arc reach principle and the lower bound}

The within--arc behaviour mirrors the hook chip--firing, localised to one block.

\begin{lemma}[Within--arc reach cost]\label{lem:reach}
Let the zero arc in block $B_k$ ($4\le k\le n-2$) have its two bounding far junctions
intact. Then the arc--scalar is nonzero only if at least $k-3$ of the arc's $k-2$ interior
positions lie in $S$; the minimum $k-3$ is attained \textup(uniquely, by a prefix of the
descending run\textup), giving a nonzero monomial.
\end{lemma}

\noindent\emph{Status.} Verified for $\la=(2,2,1,1,1)$: the $B_4$ arc (length $2$) needs
$\ge1$ insertion, the $B_5$ arc (length $3$) needs $\ge2$; minimal solutions are unique
prefixes, with nonzero value. This is the exact localisation of the hook reach principle
(a descending run of generators with rank--one boundary data); a from--scratch proof
of the cost is expected to follow the hook chip--firing argument restricted to the arc.

\begin{theorem}[Lower bound, far junctions intact]\label{thm:lb-intact}
If $S$ leaves all far junctions intact and $|S|<\tfrac{\tau(\tau+1)}2$, then $c_S=0$.
\end{theorem}

\begin{proof}
By Proposition~\ref{prop:arcs}, $c_S=\prod_j\sigma_j$. By
Proposition~\ref{prop:zeroarcs} the arcs in $B_4,\dots,B_{n-2}$ are zero at
$S=\varnothing$; for $c_S\ne0$ \emph{every} such arc must be made nonzero, and by
Lemma~\ref{lem:reach} the arc in $B_k$ requires at least $k-3$ insertions among its own
(disjoint) interior positions. Summing over the $\tau$ zero arcs,
\[
|S|\ \ge\ \sum_{k=4}^{n-2}(k-3)\ =\ \sum_{j=1}^{m-1}j\ =\ \binom m2\ =\ \frac{\tau(\tau+1)}2 . \qedhere
\]
\end{proof}

\begin{remark}[The pinch--breaking case]
A subset may instead insert a $\Pm$ at a far--junction position, breaking the pinch and
\emph{merging} two adjacent arcs. The two zero arcs of $(2,2,1,1,1)$ share the junction
$(1,5)$; breaking it merges them into one arc, which is then verified to require
$\ge3=\tfrac{\tau(\tau+1)}2$ insertions to be nonzero --- so merging does not beat the
separate cost. A general merged--arc reach bound (the analogue of the hook ``merged
junctions cost more'') would upgrade Theorem~\ref{thm:lb-intact} to the full lower bound
$\ord\ge\tfrac{\tau(\tau+1)}2$. For $(2,2,1,1,1)$ the full statement
$\min\{|S|:c_S\ne0\}=3$ is verified exhaustively over all $\binom{21}{\le3}$ subsets at
$q=5$ and $q=7/3$.
\end{remark}

\section{Upper bound and the order law}

\begin{proposition}[Critical survivors]\label{prop:surv}
For $\la=(2,2,1,1,1)$ there are exactly $14=C_4$ subsets $S$ with $|S|=3$ and $c_S\ne0$,
and every one of them has
\[
c_S=\frac{q^{8}}{(q+1)^{16}}>0 .
\]
Hence $G_3=\sum_{|S|=3}c_S=14\,q^8/(q+1)^{16}\ne0$ and, with
Theorem~\ref{thm:lb-intact} and the verified pinch--breaking bound,
\[
\ord_{s=0}P_{(2,2,1,1,1)}=3=\frac{\tau(\tau+1)}2 .
\]
\end{proposition}

\noindent The equal positive value across all critical survivors is the off--hook echo of
the hook fact that all critical $c_S$ equal $q^m/(q+1)^{2m}$; it makes the absence of
cancellation at the critical grade immediate. (Value and count verified symbolically at
$q=5$ and $q=7/3$.)

\section{Verification}

All computations are in
\texttt{\~{}/projects/scratch/2026-05-23-omega-monodromy-verify/}
(\texttt{2026-05-27-*.py}), over $\mathbb Q(q)$ evaluated at $q=5,\,7/3$, reusing the
seminormal machinery of \texttt{compute\_p\_lambda.py}.

\begin{center}
\begin{tabular}{c|c|c|c|c|c|c}
$\la$ & $m$ & $\tau$ & $\dim$ & \#zero arcs & min--support & $\ord$ pred.\\\hline
$(2,2,1,1)$   & 2 & 1 & 9  & 1 & 1 & 1\\
$(2,2,1,1,1)$ & 3 & 2 & 14 & 2 & 3 & 3\\
\end{tabular}
\end{center}

\noindent Independently confirmed for $(2,2,1^m)$: (i) $r_i\equiv4$ and the Gram--rank band
law ($\rank G_{a,b}=4$ for $|a-b|\le1$, $=1$ for $|a-b|\ge2$); (ii) the character identity
$\tfrac14(f+2\chi_2+\chi_{22})=1$ for $m=2,\dots,5$ (Murnaghan--Nakayama); (iii) the pinch
factorisation $c_\varnothing=\prod_j\sigma_j$ with exactly $\tau$ zero arcs in
$B_4,\dots,B_{n-2}$; (iv) the within--arc reach costs $k-3$ ($B_4{:}1$, $B_5{:}2$); (v) the
$14=C_4$ equal positive critical survivors for $(2,2,1,1,1)$.

\section{What is proved, and what remains}

\textbf{Proved.}
\begin{itemize}
\item The matrix--transfer factorisation $c_S=\tr\prod_r\Theta_r$
(Prop.~\ref{prop:transfer}) --- the off--hook replacement of scalar sandwiches.
\item \textbf{The rank--one pinch} (Thm.~\ref{thm:pinch}): for $|a-b|\ge2$,
$\Pq^{(a)}\Pq^{(b)}$ projects onto a joint $q$--eigenspace of dimension
$\tfrac14(f^\la+2\chi^\la_{(2)}+\chi^\la_{(2,2)})$, equal to $1$ for all $(2,2,1^m)$;
uniform because it is a class function. This is the structural heart and the genuine
generalisation of the hook rank--one collapse.
\item The pinch factorisation into arc--scalars (Prop.~\ref{prop:arcs}) and the lower
bound for all subsets that leave the far junctions intact
(Thm.~\ref{thm:lb-intact}), via $\sum_{k=4}^{n-2}(k-3)=\tfrac{\tau(\tau+1)}2$.
\end{itemize}

\textbf{Verified, not yet proved from scratch.}
\begin{itemize}
\item The within--arc reach cost $k-3$ (Lem.~\ref{lem:reach}) --- expected from a localised
hook chip--firing on the descending run with rank--one boundary.
\item The pinch--breaking (merged--arc) cost $\ge\tfrac{\tau(\tau+1)}2$, which would
complete the unconditional lower bound; verified exhaustively for $(2,2,1,1,1)$.
\item The order law $\ord=3$ for $(2,2,1,1,1)$ itself is computationally certain
(min--support $=3$ over all subsets); the contribution here is the structural proof
framework that explains it and connects it to hooks.
\end{itemize}

\paragraph{Outlook.} The rank--one pinch turns the off--hook problem into the same
reachability question as hooks, now on an \emph{arc graph} (far junctions $=$ pinch nodes,
zero arcs $=$ the blocks $B_4,\dots,B_{n-2}$). The off--hook junction cost $k-3$ replaces
the hook cost $k-2$, and the total is $\binom m2=\tfrac{\tau(\tau+1)}2$ in both worlds. The
two missing reach lemmas (within--arc and merged--arc) are the precise remaining gaps;
proving them for the family $(2,2,1^m)$ would give the order law for an infinite off--hook
family, the natural next step beyond the hook base.

\end{document}
